\documentclass[UTF8,11pt]{ctexart}
\usepackage[a4paper,margin=24mm]{geometry}
\usepackage{amsmath,amssymb,amsthm,mathtools,mathrsfs}
\usepackage[all]{xy}
\usepackage{xurl,hyperref,enumitem}
\usepackage{tikz}
\usetikzlibrary{arrows.meta,cd,decorations.markings,calc,patterns}
\hypersetup{hidelinks,breaklinks=true}
\setlength{\emergencystretch}{3em}
\allowdisplaybreaks
\newtheorem*{theorem}{Theorem}
\newtheorem*{thm}{Theorem}
\newtheorem*{lemma}{Lemma}
\newtheorem*{proposition}{Proposition}
\newtheorem*{prop}{Proposition}
\newtheorem*{corollary}{Corollary}
\newtheorem*{cor}{Corollary}
\newtheorem*{claim}{Claim}
\theoremstyle{definition}
\newtheorem*{definition}{Definition}
\newtheorem*{defn}{Definition}
\newtheorem*{remark}{Remark}
\newtheorem*{example}{Example}
\newcommand{\R}{\mathbb R}
\newcommand{\C}{\mathbb C}
\newcommand{\Z}{\mathbb Z}
\newcommand{\Q}{\mathbb Q}
\newcommand{\N}{\mathbb N}
\newcommand{\F}{\mathbb F}
\newcommand{\D}{\mathbb D}
\newcommand{\bB}{\mathbb B}
\newcommand{\bC}{\mathbb C}
\newcommand{\bR}{\mathbb R}
\newcommand{\bZ}{\mathbb Z}
\newcommand{\bN}{\mathbb N}
\newcommand{\bQ}{\mathbb Q}
\newcommand{\bD}{\mathbb D}
\newcommand{\bF}{\mathbb F}
\newcommand{\CP}{\mathbb{CP}}
\newcommand{\RP}{\mathbb{RP}}
\newcommand{\sO}{\mathscr O}
\newcommand{\sI}{\mathscr I}
\newcommand{\sF}{\mathscr F}
\newcommand{\sG}{\mathscr G}
\newcommand{\sH}{\mathscr H}
\newcommand{\sR}{\mathscr R}
\newcommand{\sM}{\mathscr M}
\newcommand{\sV}{\mathscr V}
\newcommand{\sU}{\mathscr U}
\DeclareMathOperator{\Hom}{Hom}
\DeclareMathOperator{\Ext}{Ext}
\DeclareMathOperator{\Tor}{Tor}
\DeclareMathOperator{\Spec}{Spec}
\DeclareMathOperator{\Proj}{Proj}
\DeclareMathOperator{\supp}{supp}
\DeclareMathOperator{\im}{im}
\DeclareMathOperator{\id}{id}
\DeclareMathOperator{\Tr}{Tr}
\DeclareMathOperator{\tr}{tr}
\DeclareMathOperator{\rank}{rank}
\DeclareMathOperator{\ord}{ord}
\DeclareMathOperator{\dist}{dist}
\DeclareMathOperator{\End}{End}
\DeclareMathOperator{\Aut}{Aut}
\DeclareMathOperator{\Gal}{Gal}
\DeclareMathOperator{\vol}{vol}
\DeclareMathOperator{\Cl}{Cl}
\DeclareMathOperator{\coker}{coker}
\DeclareMathOperator{\codim}{codim}
\DeclareMathOperator{\divergence}{div}
\DeclareMathOperator{\Ric}{Ric}
\DeclareMathOperator{\grad}{grad}
\DeclareMathOperator{\Hess}{Hess}
\newcommand{\abs}[1]{\left\lvert#1\right\rvert}
\newcommand{\sabs}[1]{\lvert#1\rvert}
\newcommand{\babs}[1]{\bigl\lvert#1\bigr\rvert}
\newcommand{\Babs}[1]{\Bigl\lvert#1\Bigr\rvert}
\newcommand{\bbabs}[1]{\biggl\lvert#1\biggr\rvert}
\newcommand{\BBabs}[1]{\Biggl\lvert#1\Biggr\rvert}
\newcommand{\norm}[1]{\left\lVert#1\right\rVert}
\newcommand{\snorm}[1]{\lVert#1\rVert}
\newcommand{\bnorm}[1]{\bigl\lVert#1\bigr\rVert}
\newcommand{\Bnorm}[1]{\Bigl\lVert#1\Bigr\rVert}
\newcommand{\bbnorm}[1]{\biggl\lVert#1\biggr\rVert}
\newcommand{\BBnorm}[1]{\Biggl\lVert#1\Biggr\rVert}
\newcommand{\widebar}[1]{\overline{#1}}
\newcommand{\nicefrac}[2]{\frac{#1}{#2}}
\newcommand{\linnprod}[2]{\langle#1,#2\rangle}
\newcommand{\myindex}[1]{#1}
\newcommand{\glsadd}[1]{}
\newcommand{\avoidbreak}{}
\newcommand{\sourceref}[1]{\textnormal{[原文：\nolinkurl{#1}]}}
\newcommand{\thmref}[1]{\sourceref{#1}}
\newcommand{\lemmaref}[1]{\sourceref{#1}}
\newcommand{\propref}[1]{\sourceref{#1}}
\newcommand{\corref}[1]{\sourceref{#1}}
\newcommand{\defnref}[1]{\sourceref{#1}}
\newcommand{\exampleref}[1]{\sourceref{#1}}
\newcommand{\exerciseref}[1]{\sourceref{#1}}
\newcommand{\figureref}[1]{\sourceref{#1}}
\newcommand{\sectionref}[1]{\sourceref{#1}}
\newcommand{\appendixref}[1]{\sourceref{#1}}
\newcommand{\stacksref}[2]{\href{https://stacks.math.columbia.edu/tag/#1}{#2 [#1]}}

\begin{document}
\section*{038\quad 曲线映射的Riemann--Hurwitz亏格与分歧公式}
\subsection*{来源与计数目标}
The Stacks Project Authors, \emph{The Stacks Project}。Chapter 53, Section 53.12；Lemma 53.12.2 的证明在该节前置讨论中，连同特征零分歧重数计算及 Lemma 53.12.3 收入。使用实际取得的网页数学 TeX 正文，获取日期：2026-09-28；未取得网页构建的固定提交号，不编造版本。
\par\url{https://stacks.math.columbia.edu/tag/0C1D}
\par\url{https://stacks.math.columbia.edu/tag/0C1B}
\par\url{https://stacks.math.columbia.edu/tag/0C1E}
\par\url{https://stacks.math.columbia.edu/tag/0C1A}

\par\textbf{本文件只计一个问题：}光滑适当曲线一般点 \'etale 映射的 Riemann--Hurwitz 公式，并保留原讨论中的特征零分歧项计算。
\subsection*{作者命题与人类证明}

\paragraph{53.12 Riemann--Hurwitz: the proof discussion preceding the numbered statement.}
Let $k$ be a field. Let $f:X\to Y$ be a morphism of smooth curves over $k$. Then we obtain a canonical exact sequence
\[f^*\Omega_{Y/k}\xrightarrow{\mathrm df}\Omega_{X/k}\longrightarrow\Omega_{X/Y}\longrightarrow0\]
by Morphisms, Lemma 29.33.9. Since $X$ and $Y$ are smooth, the sheaves $\Omega_{X/k}$ and $\Omega_{Y/k}$ are invertible modules, see Morphisms, Lemma 29.35.12. Assume the first map is nonzero, i.e., assume $f$ is generically \'etale, see Lemma 53.12.1. Let $R\subset X$ be the closed subscheme cut out by the different $\mathfrak D_f$ of $f$. By Discriminants, Lemma 49.12.6 this is the same as the vanishing locus of $\mathrm df$, it is an effective Cartier divisor, and we get
\[f^*\Omega_{Y/k}\otimes_{\mathcal O_X}\mathcal O_X(R)=\Omega_{X/k}.\]
In particular, if $X,Y$ are projective with $k=H^0(Y,\mathcal O_Y)=H^0(X,\mathcal O_X)$ and $X,Y$ have genus $g_X,g_Y$, then we get the Riemann--Hurwitz formula
\begin{align*}
2g_X-2&=\deg(\Omega_{X/k})\\
&=\deg\bigl(f^*\Omega_{Y/k}\otimes_{\mathcal O_X}\mathcal O_X(R)\bigr)\\
&=\deg(f)\deg(\Omega_{Y/k})+\deg(R)\\
&=\deg(f)(2g_Y-2)+\deg(R).
\end{align*}
The first and last equality by Lemma 53.8.4. The second equality by the isomorphism of invertible sheaves given above. The third equality by additivity of degrees (Varieties, Lemma 33.44.7), the formula for the degree of a pullback (Varieties, Lemma 33.44.11), and finally the formula for the degree of $\mathcal O_X(R)$ (Varieties, Lemma 33.44.9).

To use the Riemann--Hurwitz formula we need to compute $\deg(R)=\dim_k\Gamma(R,\mathcal O_R)$. By the structure of zero dimensional schemes over $k$ (see for example Varieties, Lemma 33.20.2), we see that $R$ is a finite disjoint union of spectra of Artinian local rings $R=\coprod_{x\in R}\Spec(\mathcal O_{R,x})$ with each $\mathcal O_{R,x}$ of finite dimension over $k$. Thus
\[\deg(R)=\sum_{x\in R}\dim_k\mathcal O_{R,x}=\sum_{x\in R}d_x[\kappa(x):k]\]
with
\[d_x=\operatorname{length}_{\mathcal O_{R,x}}\mathcal O_{R,x}=\operatorname{length}_{\mathcal O_{X,x}}\mathcal O_{R,x}\]
the multiplicity of $x$ in $R$ (see Algebra, Lemma 10.52.12). Let $x\in X$ be a closed point with image $y\in Y$. Looking at stalks we obtain an exact sequence
\[\Omega_{Y/k,y}\to\Omega_{X/k,x}\to\Omega_{X/Y,x}\to0.\]
Choosing local generators $\eta_x$ and $\eta_y$ of the (free rank $1$) modules $\Omega_{X/k,x}$ and $\Omega_{Y/k,y}$ we see that $\eta_y\mapsto h\eta_x$ for some nonzero $h\in\mathcal O_{X,x}$. By the exact sequence we see that $\Omega_{X/Y,x}\cong\mathcal O_{X,x}/h\mathcal O_{X,x}$ as $\mathcal O_{X,x}$-modules. Since the divisor $R$ is cut out by $h$ (see above) we have $\mathcal O_{R,x}=\mathcal O_{X,x}/h\mathcal O_{X,x}$. Thus we find the following equalities
\begin{align*}
d_x&=\operatorname{length}_{\mathcal O_{X,x}}(\mathcal O_{R,x})\\
&=\operatorname{length}_{\mathcal O_{X,x}}(\mathcal O_{X,x}/h\mathcal O_{X,x})\\
&=\operatorname{length}_{\mathcal O_{X,x}}(\Omega_{X/Y,x})\\
&=\operatorname{ord}_{\mathcal O_{X,x}}(h)\\
&=\operatorname{ord}_{\mathcal O_{X,x}}(\text{``}\eta_y/\eta_x\text{''}).
\end{align*}
The first equality by our definition of $d_x$. The second and third we saw above. The fourth equality is the definition of $\operatorname{ord}$, see Algebra, Definition 10.121.2. Note that since $\mathcal O_{X,x}$ is a discrete valuation ring, the integer $\operatorname{ord}_{\mathcal O_{X,x}}(h)$ just the valuation of $h$. The fifth equality is a mnemonic.

Here is a case where one can ``calculate'' the multiplicity $d_x$ in terms of other invariants. Namely, if $\kappa(x)$ is separable over $k$, then we may choose $\eta_x=\mathrm ds$ and $\eta_y=\mathrm dt$ where $s$ and $t$ are uniformizers in $\mathcal O_{X,x}$ and $\mathcal O_{Y,y}$ (Lemma 53.12.3). Then $t\mapsto us^{e_x}$ for some unit $u\in\mathcal O_{X,x}$ where $e_x$ is the ramification index of the extension $\mathcal O_{Y,y}\subset\mathcal O_{X,x}$. Hence we get
\[\eta_y=\mathrm dt=\mathrm d(us^{e_x})=es^{e_x-1}u\,\mathrm ds+s^{e_x}\mathrm du.\]
Writing $\mathrm du=w\,\mathrm ds$ for some $w\in\mathcal O_{X,x}$ we see that
\[\text{``}\eta_y/\eta_x\text{''}=es^{e_x-1}u+s^{e_x}w=(e_xu+sw)s^{e_x-1}.\]
We conclude that the order of vanishing of this is $e_x-1$ unless the characteristic of $\kappa(x)$ is $p>0$ and $p$ divides $e_x$ in which case the order of vanishing is $>e_x-1$.

Combining all of the above we find that if $k$ has characteristic zero, then
\[2g_X-2=(2g_Y-2)\deg(f)+\sum_{x\in X}(e_x-1)[\kappa(x):k]\]
where $e_x$ is the ramification index of $\mathcal O_{X,x}$ over $\mathcal O_{Y,f(x)}$.

\begin{lemma}[53.12.2; Tag 0C1D]
Let $f:X\to Y$ be a morphism of smooth proper curves over a field $k$ which satisfies the equivalent conditions of Lemma 53.12.1. If $k=H^0(Y,\mathcal O_Y)=H^0(X,\mathcal O_X)$ and $X$ and $Y$ have genus $g_X$ and $g_Y$, then
\[2g_X-2=(2g_Y-2)\deg(f)+\deg(R)\]
where $R\subset X$ is the effective Cartier divisor cut out by the different of $f$.
\end{lemma}
\begin{proof}
See discussion above; we used Discriminants, Lemma 49.12.6, Lemma 53.8.4, and Varieties, Lemmas 33.44.7 and 33.44.11.
\end{proof}
\begin{lemma}[53.12.3; Tag 0C1E]
Let $X\to\Spec(k)$ be smooth of relative dimension $1$ at a closed point $x\in X$. If $\kappa(x)$ is separable over $k$, then for any uniformizer $s$ in the discrete valuation ring $\mathcal O_{X,x}$ the element $\mathrm ds$ freely generates $\Omega_{X/k,x}$ over $\mathcal O_{X,x}$.
\end{lemma}
\begin{proof}
The ring $\mathcal O_{X,x}$ is a discrete valuation ring by Algebra, Lemma 10.140.3. Since $x$ is closed $\kappa(x)$ is finite over $k$. Hence if $\kappa(x)/k$ is separable, then any uniformizer $s$ maps to a nonzero element of $\Omega_{X/k,x}\otimes_{\mathcal O_{X,x}}\kappa(x)$ by Algebra, Lemma 10.140.4. Since $\Omega_{X/k,x}$ is free of rank $1$ over $\mathcal O_{X,x}$ the result follows.
\end{proof}

\subsection*{整理说明（非作者正文）}
主命题引用的 Lemma 53.12.1 条件在本文件前置作者讨论中已说明：$\mathrm df$ 非零，等价于 $f$ 一般点 \'etale，亦即函数域扩张有限可分。不是把零映射或纯不可分映射纳入该公式。标准先修为余切正合列、光滑曲线的微分层可逆、different 与微分零因子的联系、线丛次数可加和拉回次数公式、典范层次数 $2g-2$。本题承担从微分映射到分歧因子、长度与赋值，再到亏格等式的完整联系。

没有只抄 Lemma 53.12.2 的 ``See discussion above''：对应讨论及重数计算全部在前。原文随后讨论任意特征下温和分歧的充要条件并指向 Lemma 53.12.4；本题不提出该更强充要条件，故不将其未取证明混入计数。原文计算式中 $e$ 与 $e_x$ 并用，按原样保留，周围文字定义的是 $e_x$。只统一微分与长度算子排版、引号和断行；网页评论不作为证明。
\subsection*{可修改的简短标注（非作者正文）}
$c$：曲线亏格在覆盖与分歧下的变化。

$a$：Kähler 微分正合列；可逆层；different；Cartier 因子；线丛次数；模长度；离散赋值；局部参数。

\texttt{cross\_theory\_status = yes}。把映射的几何分歧编码为微分层之间映射的零因子，再用局部长度和赋值计算其次数，返回整体亏格关系。位置：节首微分正合列、$d_x$ 的长度等式及局部参数展开。

全部标注均为非穷尽、可修改初稿；未标注不表示未使用。
\subsection*{许可与修改记录}
原数学正文作者：The Stacks Project Authors。按 GNU Free Documentation License 1.2 或后续版本复用；无不变章节、封面文字或封底文字。许可全文位于 \texttt{sources/stacks/GFDL-1.2.txt}。本条整理部分沿用同一许可。2026-09-28：助手摘录、独立排版并加入来源、范围说明与初稿标注；原作者不对本次整理背书。
\end{document}
